\section{实验结果与分析}

\subsection{不同 $\eta$ 和步数下的 FID 对比}

DDIM 论文在 CIFAR-10 数据集上进行了系统性的消融实验，考察 $\eta$ 和采样步数 $S$ 对 FID 的影响。

关键观察如下：

\begin{enumerate}
    \item \textbf{全步数（$S=1000$）}：所有 $\eta$ 值的 FID 都相当，约为 3-4。DDPM（$\eta=1$）略优于确定性 DDIM（$\eta=0$）
    \item \textbf{中等步数（$S=50\sim100$）}：确定性 DDIM（$\eta=0$）显著优于 DDPM（$\eta=1$），FID 差距可达 5-10
    \item \textbf{少步数（$S=10\sim20$）}：确定性 DDIM 仍能保持较好的质量（FID $\sim$10-15），而 DDPM 的 FID 急剧恶化（$>$50）
\end{enumerate}

\begin{warningbox}{步数减少时的质量-速度权衡}
虽然确定性 DDIM 在减少步数时更为鲁棒，但步数减少到一定程度后，质量仍会不可避免地下降。这是因为较大的时间步跨度使得每步更新的``跳跃''更大，噪声预测网络在远离训练分布的区域的预测准确性会降低。DDIM 论文的实验表明，在 CIFAR-10 上约 $S=50$ 步即可达到与 1000 步 DDPM 相当的质量。
\end{warningbox}

\subsection{样本质量的定性观察}

讲者还讨论了 DDIM 生成样本的一些定性特征：

\begin{itemize}
    \item 确定性 DDIM 生成的图像通常比 DDPM 略微``锐利''（sharper），因为没有额外噪声的注入
    \item DDPM 的随机性有时反而能带来更多的多样性和``自然感''
    \item 在步数充足（$\geq$50）的情况下，两者的视觉质量差异不大
\end{itemize}

\subsection{本章小结}

实验结果验证了 DDIM 在减少采样步数方面的优势：确定性采样在步数减少时质量退化更缓慢。约 50 步的 DDIM 即可匹配 1000 步 DDPM 的质量，实现约 20 倍加速。


